\chapter[Sermon 44. The Lord Christ Performs Another, and Confirms His Elect, That They Should Not Doubt of the Faith of God's Promises, Etc.]{The Lord Christ Performs Another, and Confirms His Elect, That They Should Not Doubt of the Faith of God's Promises, Etc.}
\label{sermon:044}
\markright{Sermon 44. The Lord Christ Performs Another, and Confirms His ...}

And the angel I saw standing upon the sea and upon the earth, lifted up his hand to heaven and swore by him who lives forevermore, who created heaven and the things within it, and the sea and the things within it, that there shall be no more time; but in the days of the voice of the seventh angel, when he shall begin to blow, even the mystery of God shall be finished, as he preached through his servants the prophets.

\index{Antichrist}\index{Mahomet}But while the wicked triumphed, and the enemies of God, Antichrist and Mahomet, overcame with great success, while all good men were oppressed, and deception and lying reigned everywhere, many men will think that there shall never be an end, neither of these evils, nor yet of the world. For the Apostle Saint Peter knew this, saying that in the latter days mockers will come, who will follow their own desires, and will say, “Where is the promise of his coming?” Of whom Malachi also reasons in chapters 3 and 4. But to the intent that the goodness of God might heal the wounds of the godly, and might advance them in the truth against lying and revolt, and establish them in the same, Christ appears and swears solemnly in the sight of all men. This thing must be explained by all circumstances. For it is a matter of great weight, full of comfort, and right wholesome and necessary for all men.

\index{Daniel (Book of)}There is no doubt that he alludes to the final chapter of Daniel, wherein the Angel of the Lord swears, confirming with a solemn oath that all the things previously told to the prophet by prophecy will be fulfilled in their times. Therefore, this mighty angel swears now also, indeed, Christ himself, who set his feet on the sea and land. For by the posture and bearing of his body, he demonstrates steadfastness, lest we should doubt anything of his faith and truth; since he is lord of all, he stands moreover on feet, not fleshly, but of fiery pillars. All things therefore of Christ are certain, sure, and immovable. He who rests on him stands surely; he who believes his words shall not be confounded. And it is no new thing that Christ swears. For we read very often in Scripture that God has sworn. We read in the Gospel that the same Lord Christ has most often repeated, “Verily I say unto you,” “Verily, verily I say unto you,” which is an oath of one swearing. When Caiaphas adjured the Lord in judgment, Christ did not conceal and dissemble by holding his peace, but with express words confessed the truth. From this you may learn that the Lord, when he forbade swearing at all, did not intend the sacrament of swearing. Which, where the boorish Anabaptists will not understand, they stir up wonderful trouble, worthy to be put to silence with more severity.

\index{Hebrews, Epistle to the}But why, or to what end oaths are made or taken, the Apostle (drawing from the law in 22.2) has declared at length in the 6th chapter to the Hebrews, namely, that those who are wavering and doubtful might be confirmed, assured, and made quiet. If anyone doubts whether you are dealing faithfully with him, God commands that it be affirmed by a sacrament, so that all distrust may be removed. “Apostle,” he says, “swear by him who is greater, and to whom all controversy comes to an end, if it is confirmed by another.” In this consideration, God, intending to show more abundantly to the heirs of promise the unchangeable steadfastness of his counsel, expresses another. Even so, at this present, where divine providence foresaw that under the kingdom of Antichrist the hearts of the faithful would be most grievously tempted, and that many, because of the most prosperous fortune of Antichrist and all the wicked, would be hardened to believe God’s promises, and that many [?], which thing also Daniel in chapter 11 prophesied, would turn to Antichrist: it seemed good to God to confirm his promises by another, and that a solemn confirmation by his Son, so that those who wish to be wise may consider: if an honest and true person confirmed his promise to you by another, you would think it unworthy to doubt of his promises; how much less should it be lawful for you to doubt the promises of the Son of God, and all his words confirmed by a solemn confirmation? Believe the Son of God sworn, believe his Gospel most confirmed, even though the sky should fall and the earth gape wide. God cannot lie, which is truth, and that the eternal truth: which deceives nor is deceived: which is merciful, and loves people so that he adapts himself to their capacity. For even for us and for our infirmity he performs a sacrament, lest he should seem not to satisfy us in all things: and that all occasions of incredulity, and revolt to Antichrist, and to the filthy world might be cut away.

Now we come also to consider the manner or form of the other. Two things are here recited: the manner of the swearer, and the solemn words of the swearer. For he says, how the angel lifted up his hand toward heaven: which indeed is the most ancient rite and holy ceremony of swearers. For we read the same of Abraham in Genesis 14. And in Daniel 12, it is written of an angel: which lifting up to heaven his right hand and his left swore. We verily hold up our right hand. But where we say, that giving of voices we will hold up both our hands: we signify that we will utterly be of that sentence that we hear there propounded. Therefore the holding up of both hands doth signify a most perfect fidelity, and most assured confirmation of the thing sworn. Certainly in the holy scriptures the lifting up of the hand is often put for another. Whereof perhaps we Germans have borrowed, where we say, that is to say, thou shalt confirm me this by another. And in matters most serious and grave we are wont to use some outward ceremony, whereby we may make the words and the thing itself as it were more notable and grave. Whereupon when we pray unto God, we lift up our hands. And verily another is as it were the calling upon the name of God. Whereupon it is commonly accustomed, with great fear to perform oaths. For all men arise, and put off their caps, as they were ready to fall on their knees before the sight of God himself. When bargain or contract is made with words, the right hands are joined together also, in token of fidelity. Therefore when we take a solemn oath, we lift up our hand toward heaven, where we believe that the Lord shows himself glorious to the faithful: from whom we feel that all good things come unto us: from whence we perceive also that vengeance doth fall upon the perjured, and contemners of God. Hitherto therefore Christ applies himself unto us: and after the manner of men, to the end that men may be made the quieter, he lifts up his hands unto heaven.

And the solemn words of the swearer are these: he swears by him who lives forevermore, who made heaven and the things that are therein, and so on. Thus we read of Abraham in Genesis 14. “I lift up my hand to the Lord, God, possessor of heaven and earth.” And in Daniel 12, “He swears by him who lives forevermore.” Also in Jeremiah 4, “And you shall swear, ‘The Lord lives.’” We say so truly as God lives, and again, “So God help me.” And this is a true manner of swearing. God the creator is here most plentifully and most properly expressed, and here are all creatures severally expressed. He alone is the creator, he alone is living forevermore, as he who is life of himself and gives life to all. This creation and vivification are not communicated to others. As also he alone knows the hearts of men; that hereof we may learn to swear by the name of God alone, not to add to him any creatures, which know not the hearts, neither are life of themselves, but are less than he, and less than men, as they that are made for men. Next after God, there is nothing greater than man. Therefore let not man swear by any other than by God. For all the gentiles swear by a greater; if you swear by the saints or by the gods, you swear by men, equal verily, and not greater. God alone is greatest and best. Therefore we must swear by the name of God alone, as the scripture teaches elsewhere, in Deuteronomy 6 and 10, Exodus 23, Joshua 4 and 5, Jeremiah 45, Psalm 65, and elsewhere.

But to say that this is indeed God himself, how does he swear, you say, by him who lives forevermore—that is, by God? He swears, without doubt, by himself, as in many other passages of Scripture. Or else he swears after the dispensation and assumption of the human nature; after which he said, “My Father is greater than I,” which, notwithstanding in his deity, was nevertheless coequal with the Father.

\index{Augustine, Saint}And what I have just explained is the simplest and truest doctrine of oaths and the proper way of swearing. Yet, some understand this doctrine well enough, but for the sake of pleasing people would gladly swear by saints, and therefore ask whether they may join saints to God, especially with the reasoning, “Unless I do this, I will not be considered among the saints.” I answer that they may not, both because we have no explicit manner of swearing that we ought simply to follow in obedience to God’s honor, and also because those who require and prescribe this form would have us swear by the names of saints in heaven, and so acknowledge that we are helped and punished by their virtue and power. If you do this and acknowledge it, there is no doubt that you are gravely transgressing your sincere religion. Certainly, if you confess God before people, he will also confess you before his Father and his angels; if you deny him, he will also deny you, etc. Another [illegible] is like your confession, whereby you confess whom you acknowledge and believe to be your chief happiness, the refuge also from evil and the rewarder of good. If you therefore join saints to God himself, and equate them together, and say, “So help me God and his saints,” you are coupling them with God, granting them to be your gods, which can both help and hurt you. Therefore, take heed what you do. Read Augustine in the 145th Epistle to Publicola.

\index{Judgment, Last}However, we must also see what the angel swears by these customs and solemn words. For in this one thing consists the whole sum of the matter. The angel in the 12th chapter of Daniel swears that for a time, times, and half a time, and in the winding up, to scatter the hand of the holy people, all these things shall be fully done. So this angel here swears that there shall be no more time, but in the days of the voice of the seventh angel, when he shall begin to blow his trumpet, that the mystery of God shall be fulfilled. But here let no man understand that all time utterly, and everlastingness itself should be abolished and that there should be nothing more after the judgment; but there shall not always be such a time as now is, and such as the wicked enjoy in this world, supposing that the courses of times shall be always, that the world shall continue always, that they shall always flourish in honors and pleasures, oppressing the godly. This shall not be, says he, neither shall there be any more such a time that shall perish and be subject to changeable courses. For about the last judgment shall perish, or be renewed, all these things that shall perish, and life and glory everlasting shall succeed, I mean the time of eternity with all joy most replenished. Therefore he says not simply that there shall be no more time, but adds, in the days of the voice of the seventh angel, that is to wit, at the last judgment, that the mystery of God should be made consummate, perfect, and fully complete. What this secret, or mystery of God is, the apostle expounds and says, 1 Corinthians 15: Behold I tell you a mystery, we shall not all sleep: and the residue which are left, the mystery of God therefore is nothing else than that the end of all corruptible things is at hand, and the happy and everlasting world shall succeed; for that Christ shall then come to judgment: that Antichrist by him shall be abolished, that he with the whole body of the wicked shall be destroyed, the dead raised up again: the wicked to everlasting perdition, the godly to eternal life: for that death, sin and all corruption must be taken away from the godly, and be trodden under foot, and all misery imposed to the wicked, that they may be tormented world without end. And for as much as many times men doubt thereof, as I have said now often, Christ hath sworn that all these things shall assuredly come to pass, and that the godly shall be consummated with all glory, and that the wicked shall be consummated with all kinds of torments. And this is that great mystery of God, for which the very Son of God being incarnate, dead, and raised again from the dead ascended into Heaven, that he might convey us thither to him, having subdued Hell, Satan, Antichrist and all ungodliness. Therefore as in the 6th chapter was said to the martyrs, that they should rest for a little season, till the number of the chosen be fulfilled: so here we hear also, that the misery of God shall at length be fulfilled, \&c. This is spoken to this end also that the godly should be of quiet minds, and patiently abide deliverance. If therefore this consummation be deferred, let us abide patiently and constantly, confirmed in Christ, and his Evangelical verity, as also the apostle of Christ, Paul, hath taught us out of the Prophets, in the 10th chapter to the Hebrews.

\index{Zechariah (Book of)}Moreover, for a further declaration is added, as God has evangelized, that is, has preached a good and fortunate message, namely by his servants the prophets concerning the end of the world, the last judgment, the everlasting punishment of Antichrist and all the wicked, and the glorifying of the godly, and so forth. Neither did he say these things for a declaration only, but for confirmation also. For by the oracles of the prophets the faithful are comforted, oracles which they have never failed in anything; neither shall they deceive in the end concerning those things which they had prophesied about the last judgment. And again we see how great is the authority of the ancient scripture, and that its use is excellent in the evangelical church; wherein we see both Christ and his Apostles confirm all their sayings with prophetical scriptures, and also illuminate, set forth, and declare. The testimonies of the prophets concerning the last judgment, the reward and punishment of the godly and ungodly, the abolishing of Antichrist, of death, and of all corruption, are in Psalm 110, in Psalms 24, 26, 27, and 46, also in Daniel 7, 11, and 12, in Zechariah 14, in Malachi 3 and 4, and also elsewhere. The Apostle has cited Hosea 1 and Corinthians 15.

Therefore let us lift up our heads, brethren, let us watch and pray, for our redemption draws near. Deliver us, Christ, from all evil. Amen.
